%-------------------------------------------------------------------------------
%    ABSTRACT PAGE
%-------------------------------------------------------------------------------

\addchaptertocentry{Хураангуй} % Хураангуйг гарчигт нэмэх

\begin{center}
{\scshape\Large \univname\par} % Их сургуулийн нэр
{\scshape\large \facname\par}\vspace{0.5cm} % Сургуулийн нэр
{\huge\textbf{{Хураангуй}} \par}
\bigskip
{\Large{\ttitle} \par} % Тезисийн нэр (Artificial Intelligence Integrated digital Music processing system)
\bigskip

{\normalsize \shortname \par} % Зохиогчийн нэр
\addressname
\end{center}

\textit{\textbf{Түлхүүр үгс: Хиймэл оюун ухаан, AI Mixing, CNN, MusicVAE, LLM-based RAG, Дижитал хөгжим боловсруулалт, Хөгжмийн боловсрол}}
\bigskip

Энэхүү дипломын ажлын хүрээнд эхлэн суралцагч болон хөгжим сонирхогчдод хиймэл оюун ухаанд (AI) суурилсан хөгжмийн сургалт, боловсрол олгох нэгдсэн экосистемийг байгуулах замаар ая, хөгжим зохиох үйл явцыг хялбарчлах шийдлийг боловсрууллаа. Уг систем нь суралцагчид болон мэргэжлийн продюсеруудыг холбох нэгдсэн орон зайг бий болгосон бүтээгдэхүүн болоход чиглэгдсэн болно.

Төслийн технологийн шийдлийн хувьд дууны Spectrogram-д CNN (Convolutional Neural Networks) ашиглан анализ хийж, EQ болон компрессорын тохиргоог санал болгох "AI Mixing Assistant", хөгжмийн онол болон бүтээлч байдлын зөвлөмж өгөх "MusicVAE", мөн баталгаажсан өгөгдөлд суурилсан LLM-based RAG технологи бүхий ухаалаг чатботыг хөгжүүлсэн. 

Энэхүү ажил нь хөгжмийн "сайн" болон "муу" миксийг визуал хэлбэрээр ялгах техникийн нарийвчлалыг хангахын зэрэгцээ, анхан шатны хэрэглэгчдэд ойлгомжтой интерфэйсээр дамжуулан мэргэжлийн түвшний зөвлөгөө, сургалтын орчинг бүрдүүлснээрээ практик ач холбогдолтой юм.

